% unidades/24-forma-casual.tex
\chapter{💬 普通形 — Forma Casual}
\unidadheader{24}{普通形}{Forma Casual}{Hablar de manera informal con amigos}

\begin{tcolorbox}[phrasebox, title={⚙️ Tabla de formas: polite vs casual}]
\begin{tabular}{lllll}
\toprule
 & \textbf{Pres. +} & \textbf{Pres. -} & \textbf{Pas. +} & \textbf{Pas. -} \\
\midrule
\textbf{Polite} & 食べます & 食べません & 食べました & 食べませんでした \\
\textbf{Casual} & 食べる & 食べない & 食べた & 食べなかった \\
\midrule
\textbf{Polite} & 静かです & 静かじゃない & 静かでした & 静かじゃなかった \\
\textbf{Casual} & 静かだ & 静かじゃない & 静かだった & 静かじゃなかった \\
\midrule
\textbf{Polite} & 学生です & 学生じゃない & 学生でした & 学生じゃなかった \\
\textbf{Casual} & 学生だ & 学生じゃない & 学生だった & 学生じゃなかった \\
\bottomrule
\end{tabular}
\end{tcolorbox}

\begin{tcolorbox}[vocabbox, title={📌 Partículas y expresiones conversacionales}]
\begin{tabularx}{\linewidth}{llX}
\toprule
\textbf{Expresión} & \textbf{Uso} & \textbf{Ejemplo} \\
\midrule
〜ね & confirmar / compartir & \ruby{今日}{きょう}、\ruby{暑}{あつ}いね。 \\
〜よ & afirmar / convencer & それ、\ruby{私}{わたし}のだよ。 \\
〜よね & confirmar algo que crees & \ruby{明日}{あした}\ruby{休}{やす}みだよね？ \\
〜かな & preguntarse / dudar & いくらかな？ \\
〜でしょ & ¿ves?, ¿verdad? & ほら、\ruby{言}{い}ったでしょ。 \\
〜じゃない & ¿no es que...? & あれ、\ruby{田中}{たなか}さんじゃない？ \\
\bottomrule
\end{tabularx}
\end{tcolorbox}

\subsection*{🗣️ Frases Clave}

\frase{ねえ、\ruby{今日}{きょう}ひま？}{Oye, ¿estás libre hoy?}
\frase{うん、\ruby{暇}{ひま}だよ。}{Sí, estoy libre.}
\frase{じゃあ、いっしょに\ruby{行}{い}かない？}{Entonces, ¿no vamos juntos?}
\frase{いいね！\ruby{行}{い}こう！}{¡Genial! ¡Vamos!}

\subsection*{⚙️ Gramática}

\patron{Forma casual: verbos}
  {辞書形 (dict.) = casual presente afirmativo}
  {\ej{\ruby{明日}{あした}、\ruby{映画}{えいが}\ruby{見}{み}る？}
    {¿Mañana ves una película?}
  \ej{え、\ruby{行}{い}かないの？}{¿Eh, no vas?}
  \ej{もう\ruby{食}{た}べた？}{¿Ya comiste?}}

\patron{Invitación casual: 〜ない？}
  {V-ない (neg.) + ？ = invitación informal}
  {\ej{いっしょに\ruby{ゲーム}{ゲーム}しない？}
    {¿No jugamos juntos un videojuego?}}

\begin{tcolorbox}[ejerciciobox, title={📝 Ejercicios Rápidos}]
\begin{enumerate}
  \item Convierte a casual: \ruby{食}{た}べません → ? / \ruby{行}{い}きました → ?
  \item Escribe un mini-diálogo casual de 4 líneas entre amigos.
  \item ¿Cuándo NO debes usar la forma casual en japonés?
\end{enumerate}
\vspace{0.4em}
\textbf{Respuestas:}
1) \ruby{食}{た}べない / \ruby{行}{い}った\quad
2) (personal)\quad
3) Con personas mayores, jefes, clientes, situaciones formales — usar siempre forma polite o
keigo.
\end{tcolorbox}

\begin{tcolorbox}[kanjibox, title={🀄 Kanjis de la Unidad (repaso)}]
\begin{tabularx}{\linewidth}{cllXX}
\toprule
\textbf{Kanji} & \textbf{On} & \textbf{Kun} & \textbf{Significado} & \textbf{Vocab} \\
\midrule
\kanjirow{普}{ふ}{—}{común / ordinario}{\ruby{普通}{ふつう} / \ruby{普段}{ふだん}}
\kanjirow{通}{つう}{とお・かよ}{pasar / comunicar}{\ruby{普通}{ふつう} / \ruby{通学}{つうがく}}
\kanjirow{話}{わ}{はな・はなし}{hablar / conversación}{\ruby{話}{はなし} / \ruby{電話}{でんわ}}
\bottomrule
\end{tabularx}
\end{tcolorbox}
